\begin{table}[t]
\centering\footnotesize\setlength{\tabcolsep}{3pt}
\caption{Сертификат знаковой ошибки при нулевом пороге. Границы точны на декартовой области Q12; $U_{tr}$ и $U_{te}$ обозначают число несертифицированных решений на фиксированных обучающих и тестовых входах. Все эмпирические различия решений равны нулю. Ретроспективный анализ не меняет выбор.}
\label{tab:signed_certificate}
\begin{tabular}{rrrrrr}
\toprule
$L$ & Байт & $D_-$ & $D_+$ & $U_{tr}$ & $U_{te}$\\
\midrule
513 & 10,314 & $-10065$ & 8250 & 64 & 12\\
1025 & 20,554 & $-3426$ & 2324 & 26 & 3\\
\bottomrule
\end{tabular}
\end{table}

\paragraph{Интервал ошибки с учётом знака и совпадение пороговых решений.}
При том же переборе конечного входного домена можно сохранить знаки ошибок
отдельных рёбер. Определим $\delta_i(q)=\widetilde f_i(q)-f_i(q)$ и
\begin{equation}
 D_- = \sum_i\min_q\delta_i(q),\qquad
 D_+ = \sum_i\max_q\delta_i(q).
 \label{eq:signed_error_extrema}
\end{equation}
Для любого допустимого входа при указанных выше предположениях об арифметике
\begin{equation}
 D_-\leq\widetilde Z-Z\leq D_+,
 \qquad Z\in[\widetilde Z-D_+,\widetilde Z-D_-].
 \label{eq:signed_error_interval}
\end{equation}
Действительно, категориальные слагаемые сокращаются, а ошибка каждого
численного слагаемого лежит между найденными перебором минимумом и максимумом.
Сложение этих неравенств доказывает утверждение. Область численных входов
представляет собой декартово произведение, поэтому нижняя и верхняя границы
ошибки достигаются при независимом выборе аргумента соответствующего
экстремума для каждого ребра. Таким образом, это точные глобальные экстремумы
на данном декартовом произведении, хотя некоторые комбинации могут не
соответствовать физически возможным потокам. Для их вычисления достаточно
десяти переборов по 8,193 входа; перебор $8193^{10}$ комбинаций не требуется.

Для правила принятия решения $\mathbf1[Z\geq\tau]$ совпадение гарантируется
при выполнении одного из следующих условий; равенство порогу относится
к положительному решению:
\begin{align}
 &\widetilde Z\geq\tau\quad\text{и}\quad\widetilde Z-D_+\geq\tau,
 \quad\text{или}\nonumber\\
 &\widetilde Z<\tau\quad\text{и}\quad\widetilde Z-D_-<\tau.
 \label{eq:signed_threshold_certificate}
\end{align}
Для рассматриваемых таблиц $D_-\leq0\leq D_+$, поэтому интервал содержит
$\widetilde Z$ и условие сводится к проверке его концов.
Аналогичное рассуждение гарантирует сохранение заданного уровня риска,
если $\widetilde Z$ и оба конца лежат в одном полуоткрытом интервале между
порогами. Это прямое следствие границ ошибки, а не экспериментальная
оценка приложения с несколькими уровнями риска.

Для зафиксированных массивов с $L=1025$ получено
$(D_-,D_+)=(-3426,2324)$ (табл.~\ref{tab:signed_certificate}).
Интервал, вычисленный по оценке LUT, гарантирует совпадение 168,808 из
168,834 обучающих решений и 42,206 из 42,209 тестовых решений.
Оставшиеся 26 и 3 строки соответственно дают совпадающие решения
эмпирически, но это интервальное условие их не сертифицирует.
Уточнённый анализ не изменяет выбранную таблицу и прошивки, для которых
выполнены измерения. Сертификат вычисляется на хосте; в измеряемые ядра
на устройствах не добавлялись проверки во время исполнения или резервный
переход к коэффициентному вычислению. Найденные перебором ошибки рёбер
также совпали с результатами 81,930 независимых проверок с изменением
одной координаты в скомпилированной на хосте реализации на C.

Пусть конечная размеченная выборка содержит $N$ строк, из которых $U$
не сертифицированы. Предсказания могут различаться только на этих строках,
поэтому
$|\mathrm{Acc}_{\rm LUT}-\mathrm{Acc}_{\rm coeff}|\leq U/N$.
Для balanced accuracy, если $N_y$ обозначает число строк, а $U_y$~---
число несертифицированных строк класса $y$,
\begin{equation}
 |\mathrm{BA}_{\rm LUT}-\mathrm{BA}_{\rm coeff}|
 \leq \tfrac12\left(U_0/N_0+U_1/N_1\right).
 \label{eq:certificate_metric_bound}
\end{equation}
Для 42,209 тестовых строк $(N_0,N_1)=(10000,32209)$, а проверка интервала
с учётом знака оставляет $(U_0,U_1)=(0,3)$. Следовательно, абсолютные
изменения accuracy и balanced accuracy в худшем случае не превышают
0.00711 и 0.00466 процентного пункта соответственно. Эти границы
сравнивают два зафиксированных целочисленных представления на данной
выборке; они не гарантируют правильность относительно истинной разметки
или качество на новом распределении. Наблюдаемые изменения равны нулю.

\paragraph{Проверка экспорта до выпуска прошивки.}
Сертификат можно использовать для проверки экспорта перед выпуском прошивки.
Запись о выпуске может связывать хеши коэффициентов и LUT, кодирование
признаков, правила арифметики и порог решения с вычисленным перебором
интервалом ошибки. Этот интервал определяет, для каких оценок LUT
совпадение с эталоном сертифицировано. На явно заданной выборке для проверки
та же процедура выявляет несертифицированные строки для непосредственного
сравнения с коэффициентной реализацией. Так граница ошибки на всей области
отделяется от совпадения, проверенного только на наблюдаемых входах.
В будущем автоматизированная сборка могла бы отклонять экспорт, если число
несертифицированных решений превышает допустимый предел, заданный по
обучающим или валидационным входам до обращения к тесту. Это возможное
применение анализа, а не реализованная проверка в CI или правило выбора
в настоящем исследовании. Предел на фиксированной выборке не ограничивает
частоту несертифицированных решений в новой сети. Для обработки
неопределённых оценок на устройстве потребовались бы отдельная реализация
проверки и заданного резервного вычисления, а также измерения их памяти
и времени выполнения. Настоящая предварительная проверка не добавляет
такой защиты в измеренную прошивку.

\paragraph{Ретроспективное сравнение памяти и неопределённости.}
Таблица с $L=513$ занимает 10,314 байт и также не имеет наблюдаемых
расхождений решений на обучающей и тестовой выборках. Её интервал ошибки
с учётом знака равен $[-10065,8250]$, а интервал, вычисленный по оценке LUT,
оставляет несертифицированными 64 обучающих и 12 тестовых решений.
Переход к $L=1025$ добавляет 10,240 байт параметров (рост в
1.993 раза) и уменьшает эти количества до 26 и 3. При исходном
симметричном условии на отступ эталонной оценки соответствующие
количества уменьшаются с 95 до 34 на обучающей выборке и с 21 до 4
на тестовой. Для этих двух зафиксированных таблиц дополнительная память
сужает область неопределённости, но не обеспечивает полную гарантию
сохранения решений и не устанавливает минимальный объём хранения
для заданного допустимого уровня неопределённости.
Анализ тестовой выборки для $L=513$ является ретроспективным и не
используется для изменения выбора по обучающим данным или
зафиксированного развёртывания.
